\thispagestyle{empty}
\section*{Kurzdarstellung}\label{sec:kurzdarstellung}

Entwicklungsteams setzen zunehmend auf Microservice-Architekturen, um die Flexibilität und Skalierbarkeit ihrer Anwendungen zu erhöhen.
Dabei ist aufgrund komplexerer Entwicklungsprozesse und kürzerer Release-Zyklen die Sicherstellung der Softwarequalität durch regelmäßige Softwaretest eine zentrale Herausforderung.
Klassische Testmethoden wie Integrationstests stoßen hierbei an ihre Grenzen, da sie oft zeitintensiv und fehleranfällig sind.
Diese Arbeit hat das Ziel die Eignung von Contract Testing zur Verbesserung der Schnittstelleninteroperabilität, Testeffizienz und Fehlererkennung systematisch zu untersuchen.

Zunächst erfolgte eine vergleichende Bewertung von fünf relevanten Contract-\-Testing-\-Frame\-works (\texttt{Pact}, \texttt{Spring Cloud Contract}, \texttt{Specmatic}, \texttt{Karate}, \texttt{ContractCase}) anhand definierter technischer, organisatorischer und funktionaler Kriterien.
Im Anschluss wurde das Framework \texttt{Pact} für die prototypische Implementierung von Contract Testing ausgewählt und in der realen Microservice-Anwendung \textit{msg.Sensorik} umgesetzt.
Die Implementierung wurde hinsichtlich ihrer Fähigkeit zur Fehlererkennung, ihrer Effizienz im Vergleich zu klassischen Integrationstests sowie ihres Beitrags zur Reduktion der Fehlerkosten evaluiert.
Die Ergebnisse zeigen, dass Contract Testing als ergänzende Testmethode signifikant zur Verbesserung der Softwarequalität und zur frühzeitigen Fehlererkennung beitragen kann, insbesondere in CI/CD-gestützten Entwicklungsprozessen.

\section*{Abstract}\label{sec:abstract}

